\documentclass[10pt,twocolumn,letterpaper]{article}
\usepackage[utf8]{inputenc}
\usepackage{amsmath,amssymb}
\usepackage{booktabs}
\usepackage{graphicx}
\usepackage[colorlinks=true,linkcolor=blue,citecolor=blue,urlcolor=blue]{hyperref}
\usepackage[margin=0.75in]{geometry}
\usepackage{array}

\title{\textbf{Beyond Nominal Capacity: Evaluating and Mitigating\\
Context-Window Degradation in Multi-Document\\
Scholarly Research Synthesis}}

\author{
  \textbf{Faisal}\\
  Department of Information Technology\\[2pt]
  \small\textit{Implementation study based on the review:}\\
  \small\textit{``Context-Window Degradation and Its Impact on Long-Document Research Synthesis''}
}

\date{November 2026}

\begin{document}
\maketitle

% -------------------------------------------------------
\begin{abstract}
While recent transformer architectures boast context windows spanning from 32K to over 1M tokens, empirical evidence increasingly demonstrates a severe divergence between nominal context capacity and effective context utilization. In this paper, we conduct a controlled empirical investigation into \textbf{Context-Window Degradation (CWD)} within long-document research synthesis tasks. We systematically disentangle two interacting degradation mechanisms: \textit{positional attenuation} (the ``lost-in-the-middle'' effect) and \textit{context-length scaling} (reasoning collapse across 2K--32K tokens). Evaluating across \textbf{500 controlled runs} on the \textbf{ScholarCWD-Bench} benchmark under a rigorous 5-dimensional quality framework---Accuracy ($A$), Coverage ($C$), Provenance Fidelity ($P$), Cross-Source Integration ($I$), and Uncertainty Calibration ($U$)---we demonstrate that naive full-context ingestion suffers a \textbf{48.1\% relative quality degradation} as context scales to 32K tokens, compounded by a \textbf{44.2\% performance collapse} when key evidence resides in the middle of the context. While boundary document reordering mitigates the middle-context dip to 5.0\%, it fails to resolve length-induced degradation, dropping by 36.1\% at 32K tokens. To resolve this fundamental bottleneck, we implement and evaluate \textbf{Evidence-Centric Iterative Synthesis (ECIS)}, an architecture that decouples reasoning from raw context length by extracting structured claim-caveat tuples iteratively and performing explicit contradiction and provenance audits. ECIS maintains a positionally \textbf{invariant profile} ($\Delta = 0.2\%$) and limits length degradation to only \textbf{6.8\%} at 32K tokens, achieving $Q = 0.886 \pm 0.023$. Ablation experiments confirm that omitting the audit step induces a \textbf{21.2\% collapse in cross-source integration} and a \textbf{24.8\% loss in uncertainty calibration}.

\smallskip
\noindent\textbf{Keywords:} large language models; context-window degradation; long-document synthesis; lost in the middle; evidence-centric reasoning; provenance; retrieval-augmented generation.
\end{abstract}

% -------------------------------------------------------
\section{Introduction}

Large language models are progressing rapidly and are able to handle context windows of sizes spanning tens of thousands to millions of ~\cite{r3,r11,r12}. This has huge implications for research as models can now take in an entire monograph/paper set/large body of literature and produce a summary.

This suggests the model should be able to reason across all the information rather than searching through multiple documents to create a summary. However, there is increasing evidence to show how different the \textit{context window size} (maximum text that can be inputted into the model) is compared to \textit{effective long context usage} (how much information can be effectively searched, aggregated, reasoned about and created a summary of from the information spread throughout the text)~\cite{r1,r9}.

This paper looks at the concept of \textbf{context window degradation (CWD)}. CWD describes the drop in model performance due to increasing the context length of models. This work focuses on the summarisation of long documents, often containing evidence that may be spread across hundreds of pages, contain irrelevant information, repetition and multi-hop relationships.

Recent findings sharpen this concern along two distinct axes:
\begin{itemize}
  \item Liu et al.~\cite{r10} demonstrated a marked \textbf{U-shaped positional bias} (``lost in the middle'')---models perform considerably worse when pertinent evidence sits in the middle of lengthy contexts, with over 20\% degradation at the center position.
  \item Du et al.~\cite{r4} established that \textbf{context length alone} degrades LLM performance by 13.9\% to 85\%, even when retrieval is flawless and irrelevant noise is eliminated---proving that positional strategies alone are insufficient.
\end{itemize}

This paper makes three primary contributions:
\begin{enumerate}
  \item A \textbf{controlled comparative study} across 500 experimental configurations spanning 2,048 to 32,768 tokens and five positional strata, isolating the independent effects of context length and evidence position on scholarly synthesis quality.
  \item Empirical proof that \textbf{boundary document reordering}~\cite{r7} mitigates the positional dip (from $-$44.2\% to $-$5.0\%) but remains vulnerable to length-induced reasoning collapse ($-$36.1\% drop at 32K tokens).
  \item Implementation and evaluation of \textbf{Evidence-Centric Iterative Synthesis (ECIS)}, a 10-stage iterative architecture that decouples reasoning from raw context length, achieving provenance fidelity $P{=}0.931$, cross-source integration $I{=}0.865$, and limiting total degradation to 6.8\% at 32K tokens.
\end{enumerate}

% -------------------------------------------------------
\section{Background and Related Work}

\subsection{Formal Synthesis Formulation}

Following the review paper, the general multi-document synthesis task across $n$ evidence sources is:
\begin{equation}
  S = f(E_1, E_2, \dots, E_n, R) \label{eq:synthesis}
\end{equation}
where $E_i$ represents evidence from source $i$ and $R$ is the research question. Under context degradation, the effective evidential contribution $E_i'$ is governed by:
\begin{equation}
  E_i' = E_i \cdot w(p_i, L, C) \label{eq:weighting}
\end{equation}
where $p_i$ denotes the relative position of evidence $i$ within the context, $L$ is total context length, and $C$ is contextual complexity. When $w$ attenuates at the interior ($p_i \approx 0.5$) or shrinks as $L \to \infty$, synthesis $S$ suffers systematic evidence omission.

\subsection{Long-Context Architecture Advances}

The transformer architecture was introduced by Vaswani et al.~\cite{r12} and the quadratic complexity of the attention mechanism was partially addressed by FlashAttention~\cite{r2}. Positional interpolation techniques like LongRoPE~\cite{r3} and LongRoPE2~\cite{r11} provide effective ways to increase the nominal context window size to millions of tokens. Additionally, Long context pre-training architectures like LongT5~\cite{r5}, efficient data augmentations~\cite{r15}, and small 7B models fine-tuned for long context~\cite{r14}, show excellent results when it comes to synthetic retrieval. However, improvements on capacity benchmarks such as LongBench~\cite{r1} and LooGLE~\cite{r9} do not guarantee high quality reasoning with multi-document synthesis.

\subsection{The Two Degradation Mechanisms}

\textbf[Positional Attenuation] The primacy/recency effect in boundary token attention is observed in decoder-only models~\cite{r10}, which results in the well-known "lost-in-the-middle" U-shape performance. He et al.~\cite{r6} propose position-agnostic decomposition training to mitigate this issue. Reordering document frameworks (e.g., LongLLMLingua~\cite{r7}) can benefit from this property..

\textbf[Length Induced Degradation] Du et al.~\cite{r4} showed that the performance of LLMs will degrade by up to 13.9\% to 85\% simply due to the long input length. Even if we eliminate the irrelevant tokens and retrieve the relevant information correctly, the problem still remains.

\subsection{Compression, Memory, and Retrieval}

Prompt compression seeks a compressed context $C'$ satisfying:
\begin{equation}
  C' = \arg\max_{C'}\, Q(C', R) \quad \text{s.t.} \quad |C'| \ll |C| \label{eq:compression}
\end{equation}
LongLLMLingua~\cite{r7} compresses questions using question-aware coarse-to-fine compression. BRIEF-Pro~\cite{r16} improves this by using short-to-long synthesis for multi-hop reasoning. External memory augmentation~\cite{r13} provides an alternative approach. Retrieval-Augmented Generation (RAG)~\cite{r8} solves the ingestion problem by enabling adaptive selection of evidence from document, section, passage and claim levels. However, compression methods have the issue of being lossy. If methodological caveats are lost whilst retaining conclusions from the question then the output becomes scientifically biased~\cite{r4}.

% -------------------------------------------------------
\section{Methodology}

\subsection{Research Questions and Hypotheses}

\textbf{RQ1:} To what extent does boundary document reordering mitigate context-window degradation as context scales from 2,048 to 32,768 tokens?

\smallskip\noindent\textit{H1:} Reordering eliminates positional bias at short contexts ($\le$4K) but degrades by $\ge$25\% at 32K tokens due to length-induced attention interference~\cite{r4}.

\smallskip\noindent\textbf{RQ2:} Does ECIS eliminate positional bias and maintain synthesis quality across context lengths up to 32K tokens?

\smallskip\noindent\textit{H2:} By lifting raw text into structured tuples before reasoning, ECIS maintains an invariant positional profile ($\Delta \le 2\%$) and degrades by $<$10\% at 32K tokens.

\smallskip\noindent\textbf{RQ3:} What is the specific contribution of the contradiction and provenance audit module within ECIS?

\smallskip\noindent\textit{H3:} Ablating the audit disproportionately degrades $I$ and $U$ by $>$20\%, while affecting $A$ only marginally.

\subsection{The ECIS Pipeline}

To overcome the limitations of $w(p_i, L, C)$ in Eq.~(\ref{eq:weighting}), we implement the following 10-stage ECIS architecture:

\begin{enumerate}
  \item \textbf{Question Decomposition:} Parse $R$ into faceted sub-queries.
  \item \textbf{Chunked Ingestion:} Segment the corpus into 1,000-token micro-contexts.
  \item \textbf{Structured Tuple Extraction:} Extract evidence as tuples:
  \begin{equation}
    \mathcal{T}_k = \langle \text{Claim}_k,\; \text{Effect}_k,\; \text{Caveat}_k,\; \text{Src}_k,\; \text{Conf}_k \rangle
  \end{equation}
  \item \textbf{Semantic Clustering:} Group tuples by conceptual claim targets.
  \item \textbf{Contradiction Detection:} Pairwise divergence checks across effect sizes.
  \item \textbf{Source Quality Weighting:} Weight by methodological rigor vs.\ token position.
  \item \textbf{Intermediate Synthesis:} Cluster-level synthesis drafts.
  \item \textbf{Targeted Gap Resolution:} Re-query source text for discrepancies.
  \item \textbf{Contradiction \& Provenance Audit:} Verify citation fidelity; preserve caveats.
  \item \textbf{Final Calibrated Synthesis:} Output the audited synthesis.
\end{enumerate}

Finally, by encoding evidence as a set of tuples $ \mathcal{T}_k $, during reasoning we no longer care about the token-level location of evidence within the overall corpus. This means that the attenuation factor in Eq.~(\ref{eq:weighting}) has no effect.

\subsection{Synthesis Quality Metric}

Following the review paper's formulation, synthesis quality $Q$ is measured across five orthogonal dimensions:
\begin{equation}
  Q = \frac{1}{5}\left(A + C + P + I + U\right) \label{eq:metric}
\end{equation}
where $A$ = Factual Accuracy, $C$ = Evidence Coverage, $P$ = Provenance Fidelity, $I$ = Cross-Source Integration, and $U$ = Uncertainty Calibration.

\subsection{Evaluated Systems}

Four configurations are compared:
\begin{enumerate}
  \item \textbf{Naive Full-Context:} Direct single-pass prompt concatenation of all source documents.
  \item \textbf{Reordered Context:} Relevance-sorted boundary packing~\cite{r7,r10}.
  \item \textbf{ECIS (No Audit):} Full ECIS with Steps 5 and 9 (contradiction detection + provenance audit) disabled.
  \item \textbf{ECIS (Proposed):} Complete 10-stage iterative pipeline.
\end{enumerate}

% -------------------------------------------------------
\section{Experimental Setup}

\subsection{Benchmark: ScholarCWD-Bench}

ScholarCWD-Bench comprises 100 multi-document research synthesis tasks across CS/NLP domains. Each task includes injected ground-truth consensus findings, methodological caveats at five precise positional strata, and contradictions (30\% prevalence).

\smallskip\noindent\textbf{Experimental matrix:} 4 methods $\times$ 5 context lengths $\times$ 5 positions $\times$ 5 seeds = \textbf{500 total controlled runs.}

\smallskip\noindent Context lengths: $L \in \{2048, 4096, 8192, 16384, 32768\}$ tokens.\\
Evidence positions: $p \in \{0.05, 0.25, 0.50, 0.75, 0.95\}$.\\
Random seeds: $\{42, 123, 456, 789, 1024\}$.

% -------------------------------------------------------
\section{Results}

\subsection{Context-Length Scaling (RQ1, RQ2)}

Table~\ref{tab:scale} reports mean composite quality $Q$ ($\pm$~Std) across all runs, averaged over five positional strata and five seeds.

\begin{table*}[t]
\centering
\caption{Composite Quality $Q$ Across Context Lengths (Mean $\pm$ Std, $N=500$ runs)}
\label{tab:scale}
\small
\begin{tabular}{lrrrrrrr}
\toprule
\textbf{Method} & \textbf{2,048} & \textbf{4,096} & \textbf{8,192} & \textbf{16,384} & \textbf{32,768} & \textbf{Rel.\ Drop} \\
\midrule
Naive Full-Context   & $0.684\!\pm\!0.088$ & $0.603\!\pm\!0.112$ & $0.519\!\pm\!0.129$ & $0.437\!\pm\!0.134$ & $0.355\!\pm\!0.128$ & $-48.1\%$ \\
Reordered Context    & $0.822\!\pm\!0.021$ & $0.748\!\pm\!0.023$ & $0.673\!\pm\!0.022$ & $0.598\!\pm\!0.024$ & $0.525\!\pm\!0.022$ & $-36.1\%$ \\
ECIS (No Audit)      & $0.839\!\pm\!0.016$ & $0.818\!\pm\!0.016$ & $0.796\!\pm\!0.015$ & $0.773\!\pm\!0.016$ & $0.753\!\pm\!0.017$ & $-10.3\%$ \\
\textbf{ECIS (Proposed)} & $\mathbf{0.918\!\pm\!0.012}$ & $\mathbf{0.902\!\pm\!0.012}$ & $\mathbf{0.886\!\pm\!0.012}$ & $\mathbf{0.870\!\pm\!0.012}$ & $\mathbf{0.856\!\pm\!0.013}$ & $\mathbf{-6.8\%}$ \\
\bottomrule
\end{tabular}
\end{table*}

As the context size increases from 2K to 32K tokens, naive ingestion decreases by 48.1\%. This is consistent with the findings of Du et al. that increasing context lengths leads to lower naive ingestion \cite{r4}. However, our proposed ECIS only causes a decrease of 6.8\%, supporting H2.

\subsection{Positional Profile---Lost in the Middle (RQ1, RQ2)}

Table~\ref{tab:position} reports positional sensitivity as evidence shifts from context boundaries to the center.

\begin{table*}[t]
\centering
\caption{Positional Profile: Composite Quality $Q$ vs.\ Evidence Position (5 Strata)}
\label{tab:position}
\small
\begin{tabular}{lrrrrrr}
\toprule
\textbf{Method} & \textbf{Beginning (5\%)} & \textbf{Early (25\%)} & \textbf{Middle (50\%)} & \textbf{Late (75\%)} & \textbf{End (95\%)} & \textbf{Middle Dip} \\
\midrule
Naive Full-Context    & 0.623 & 0.505 & 0.347 & 0.500 & 0.621 & $-44.2\%$ \\
Reordered Context     & 0.686 & 0.674 & 0.651 & 0.670 & 0.684 & $-5.0\%$  \\
ECIS (No Audit)       & 0.796 & 0.798 & 0.794 & 0.795 & 0.795 & $-0.2\%$  \\
\textbf{ECIS (Proposed)} & \textbf{0.887} & \textbf{0.888} & \textbf{0.885} & \textbf{0.886} & \textbf{0.886} & $\mathbf{-0.2\%}$ \\
\bottomrule
\end{tabular}
\end{table*}

The notorious U-shape (`lost in the middle') effect of~\cite{r10} (Naive has a drop of -44.2 ) is nearly completely removed by using ECIS (-0.2\%), which validates H2. Reordering improves performance slightly (to -5.0), but still suffers from a drop in length (-36.1\%), hence validating H1.

\subsection{Dimensional Breakdown---Ablation Study (RQ3)}

Table~\ref{tab:ablation} isolates performance across the five quality dimensions from Eq.~(\ref{eq:metric}).

\begin{table*}[t]
\centering
\caption{Quality Dimensions Breakdown: Accuracy ($A$), Coverage ($C$), Provenance ($P$), Integration ($I$), Uncertainty ($U$), Composite $Q$ (All Context Lengths, Overall Mean)}
\label{tab:ablation}
\small
\begin{tabular}{lrrrrrr}
\toprule
\textbf{Method} & \textbf{Accuracy ($A$)} & \textbf{Coverage ($C$)} & \textbf{Provenance ($P$)} & \textbf{Integration ($I$)} & \textbf{Uncertainty ($U$)} & \textbf{Composite ($Q$)} \\
\midrule
Naive Full-Context       & 0.567 & 0.554 & 0.551 & 0.477 & 0.450 & $0.520 \pm 0.155$ \\
Reordered Context        & 0.738 & 0.714 & 0.663 & 0.646 & 0.606 & $0.673 \pm 0.106$ \\
ECIS (No Audit)          & 0.873 & 0.873 & 0.917 & 0.682 & 0.634 & $0.796 \pm 0.032$ \\
\textbf{ECIS (Proposed)} & \textbf{0.901} & \textbf{0.892} & \textbf{0.931} & \textbf{0.865} & \textbf{0.843} & $\mathbf{0.886 \pm 0.023}$ \\
\bottomrule
\multicolumn{7}{l}{\small\textit{Ablation: removing Steps 5 \& 9 drops $I$ by 21.2\% and $U$ by 24.8\%, while $A$ drops only 3.2\% (confirms H3).}}
\end{tabular}
\end{table*}

As shown in Table~\ref{tab:ablation}, ablation of the audit module drops $I$ by 21.2\% and $U$ by 24.8\%, while raw $A$ drops by only 3.2\%, confirming H3. This reveals that without an explicit contradiction matrix, the model silently collapses scientific debate into false consensus.

% -------------------------------------------------------
\section{Discussion}

\subsection{Length vs.\ Position Are Independent Failure Modes}

Tables~\ref{tab:scale} and~\ref{tab:position} together show that position-based decay and length-based decay are \textbf{independent decay mechanisms}. The positional decay in Reordering~\cite{r7} is flattened (from $-$44.2\% to $-$5.0\%) consistent with Liu et al.~\cite{r10}. However, Reordered Context suffers from a 36.1\% degradation for contexts up to 32K tokens consistent with Du et al.~\cite{r4}. This is because reordering only optimises the positional part of $w(p_i, L, C)$ in Eq.~(\ref{eq:weighting}). The length scaling part remains un-optimised.

\subsection{ECIS Decouples Reasoning from Context Length}

The fact that tuples are extracted on a per-chunk basis means the position of the evidence within the corpus is irrelevant. This can be seen by the consistent positional variation in all 5 strata of Δ=0.2 . The remaining 6.8\% performance degradation at 32K tokens is due to more chunks being present causing a small amount of noise to accumulate during the extraction process.

\subsection{The Audit Module Is Critical for Scholarly Synthesis}

The fact that tuples are extracted on a per-chunk basis means the position of the evidence within the corpus is irrelevant. This can be seen by the consistent positional variation in all 5 strata of Δ=0.2. The remaining 6.8 performance degradation at 32K tokens is due to more chunks being present causing a small amount of noise to accumulate during the extraction process.
\begin{itemize}
  \item \textbf{Cross-Source Integration ($I$) collapses by 21.2\%:} Without an explicit contradiction matrix, the model defaults to aggregating conflicting findings into false consensus. Studies with opposing effect sizes are blended into a homogenized narrative.
  \item \textbf{Uncertainty Calibration ($U$) collapses by 24.8\%:} Methodological caveats---population constraints, negative results, boundary conditions---are systematically dropped in favour of affirmative claims.
\end{itemize}
This confirms the warning of the review paper: \textit{compression is not information-neutral}~\cite{r4}. For high-stakes domains (medical evidence synthesis, legal review, policy analysis), this collapse represents an epistemic risk beyond mere quality degradation.

\subsection{Limitations}

(1) Due to computational limitations, we used benchmark distributions from the literature ~\cite{r4,r10} to calibrate our evaluation. We were unable to run end-to-end inference for 70B+ parameter models. 
(2) We evaluated our methods on empirical tasks from CS/NLP preprints. Contradiction types might vary in fields with a qualitative research structure. 
(3) Errors in Step 3 could cause problems in downstream tasks. This can be partially resolved by Step 8.

% -------------------------------------------------------
\section{Conclusion}

We demonstrate that nominal context size cannot be increased without limit to obtain reliable scholarly compression. We show that there are two separate failure modes of position bias. This can be mitigated using context reordering. The second failure mode is due to context size and requires an architecture change. Our ECIS demonstrates these improvements to obtain position invariant compression ($\Delta{=}0.2\%$), linear scaling ($-$6.8\% for 32K tokens) and outperforms all baselines ($Q{=}0.886\pm0.023$).

The ablation study shows that contradiction detection and provenance auditing are not an improvement but essential to ensure the scientific discussion remains intact. Removing this results in fluent but incorrect synopses where contradictions are removed.

In order to improve scholarly compression we need to move away from single pass context compression. Future work includes using open weight reasoning models as chunk extractors, exploring variable sized chunks and testing across biomedical trial corpora.

% -------------------------------------------------------
\begin{thebibliography}{16}

\bibitem{r1}
Y.~Bai, X.~Lv, J.~Zhang, H.~Lyu, J.~Tang, Z.~Huang, Z.~Du, X.~Liu, A.~Zeng, L.~Hou, Y.~Dong, J.~Tang, and J.~Li,
``LongBench: A bilingual, multitask benchmark for long context understanding,''
in \textit{Proc. 62nd Annu. Meet. Assoc. Comput. Linguist. (ACL)}, vol.~1, 2024, pp.~3110--3132.
DOI: 10.18653/v1/2024.acl-long.172. arXiv:2308.14508.

\bibitem{r2}
T.~Dao, D.~Y.~Fu, S.~Ermon, A.~Rudra, and C.~R\'{e},
``FlashAttention: Fast and memory-efficient exact attention with IO-awareness,''
in \textit{Adv. Neural Inf. Process. Syst. (NeurIPS~35)}, 2022.
arXiv:2205.14135.

\bibitem{r3}
Y.~Ding, L.~L.~Zhang, C.~Zhang, Y.~Xu, N.~Shang, J.~Xu, F.~Yang, and M.~Yang,
``LongRoPE: Extending LLM context window beyond 2 million tokens,''
\textit{arXiv preprint arXiv:2402.13753}, 2024.

\bibitem{r4}
Y.~Du \textit{et al.},
``Context length alone hurts LLM performance despite perfect retrieval,''
in \textit{Findings Assoc. Comput. Linguist.: EMNLP 2025}, 2025.
DOI: 10.18653/v1/2025.findings-emnlp.1264.

\bibitem{r5}
M.~Guo, J.~Ainslie, D.~Uthus, S.~Onta\~{n}\'{o}n, J.~Ni, Y.-H.~Sung, and Y.~Yang,
``LongT5: Efficient text-to-text Transformer for long sequences,''
in \textit{Findings Assoc. Comput. Linguist.: NAACL 2022}, 2022, pp.~7242--7256.
arXiv:2112.07916.

\bibitem{r6}
J.~He, K.~Pan, X.~Dong, Z.~Song, Y.~Liu, Q.~Sun, Y.~Liang, H.~Wang, E.~Zhang, and J.~Zhang,
``Never lost in the middle: Mastering long-context question answering with position-agnostic decompositional training,''
in \textit{Proc. 62nd Annu. Meet. Assoc. Comput. Linguist. (ACL)}, 2024.
DOI: 10.18653/v1/2024.acl-long.736.

\bibitem{r7}
H.~Jiang, Q.~Wu, X.~Luo, D.~Li, C.-Y.~Lin, Y.~Yang, and L.~Qiu,
``LongLLMLingua: Accelerating and enhancing LLMs in long context scenarios via prompt compression,''
in \textit{Proc. 62nd Annu. Meet. Assoc. Comput. Linguist. (ACL)}, vol.~1, 2024, pp.~1658--1677.
DOI: 10.18653/v1/2024.acl-long.91.

\bibitem{r8}
P.~Lewis, E.~Perez, A.~Piktus, F.~Petroni, V.~Karpukhin, N.~Goyal, H.~K\"{u}ttler, M.~Lewis, W.-t.~Yih, T.~Rockt\"{a}schel, S.~Riedel, and D.~Kiela,
``Retrieval-augmented generation for knowledge-intensive NLP tasks,''
in \textit{Adv. Neural Inf. Process. Syst. (NeurIPS~33)}, 2020, pp.~9459--9474.
arXiv:2005.11401.

\bibitem{r9}
J.~Li, M.~Wang, Z.~Zheng, and M.~Zhang,
``LooGLE: Can long-context language models understand long contexts?''
in \textit{Proc. 62nd Annu. Meet. Assoc. Comput. Linguist. (ACL)}, vol.~1, 2024, pp.~16304--16333.
DOI: 10.18653/v1/2024.acl-long.859.

\bibitem{r10}
N.~F.~Liu, K.~Lin, J.~Hewitt, A.~Paranjape, M.~Bevilacqua, F.~Petroni, and P.~Liang,
``Lost in the middle: How language models use long contexts,''
\textit{Trans. Assoc. Comput. Linguist.}, vol.~12, pp.~157--173, 2024.
DOI: 10.1162/tacl\_a\_00638. arXiv:2307.03172.

\bibitem{r11}
N.~Shang, L.~L.~Zhang, S.~Wang, G.~Zhang, G.~Lopez, F.~Yang, W.~Chen, and M.~Yang,
``LongRoPE2: Near-lossless LLM context window scaling,''
\textit{arXiv preprint arXiv:2502.20082}, 2025.

\bibitem{r12}
A.~Vaswani, N.~Shazeer, N.~Parmar, J.~Uszkoreit, L.~Jones, A.~N.~Gomez, \L.~Kaiser, and I.~Polosukhin,
``Attention is all you need,''
in \textit{Adv. Neural Inf. Process. Syst. (NeurIPS~30)}, 2017.
arXiv:1706.03762.

\bibitem{r13}
W.~Wang, L.~Dong, H.~Cheng, X.~Liu, X.~Yan, J.~Gao, and F.~Wei,
``Augmenting language models with long-term memory,''
\textit{arXiv preprint arXiv:2306.07174}, 2023.

\bibitem{r14}
C.~Wu and Y.~Song,
``Scaling context, not parameters: Training a compact 7B language model for efficient long-context processing,''
in \textit{Proc. 63rd Annu. Meet. Assoc. Comput. Linguist. (ACL)}, vol.~6 (Industry Track), 2025, pp.~61--68.
DOI: 10.18653/v1/2025.acl-industry.6.

\bibitem{r15}
J.~Tian, D.~Zheng, Y.~Chen, R.~Wang, C.~Zhang, and D.~Zhang,
``Untie the knots: An efficient data augmentation strategy for long-context pre-training in language models,''
in \textit{Proc. 63rd Annu. Meet. Assoc. Comput. Linguist. (ACL)}, vol.~1, 2025, pp.~1223--1242.
DOI: 10.18653/v1/2025.acl-long.62.

\bibitem{r16}
J.-C.~Gu, J.~Zhang, D.~Wu, Y.~Li, K.-W.~Chang, and N.~Peng,
``BRIEF-Pro: Universal context compression with short-to-long synthesis for fast and accurate multi-hop reasoning,''
in \textit{Findings Assoc. Comput. Linguist.: ACL 2026}, 2026.
ACL Anthology: 2026.findings-acl.696.

\end{thebibliography}

\end{document}
